\chapter{Zasady redagowania tekstu}
\label{chap:redakcja}

Rozdział pokazuje, jak w~tej formatce wstawia się wzory, rysunki, tabele
i~fragmenty kodu. Kod źródłowy rozdziału (plik
\texttt{rozdzialy/03\_redakcja.tex}) można traktować jako ściągę --
wystarczy skopiować potrzebny fragment i~podmienić treść.

\section{Styl wypowiedzi}
\label{sec:styl}

Pracę pisze się w~formie bezosobowej. Unika się pierwszej osoby liczby
pojedynczej i~mnogiej, a~także zwrotów potocznych. Porównanie wariantów
poprawnych i~niepoprawnych zawiera tabela~\ref{tab:styl}.

\begin{table}[htbp]
  \centering
  \caption{Przykłady formy bezosobowej}
  \label{tab:styl}
  \begin{tabularx}{\linewidth}{@{}Xc@{}}
    \toprule
    \textbf{Zdanie} & \textbf{Ocena} \\
    \midrule
    Wyniki obliczeń strat ciśnienia zestawiono w~tabeli~3.4. & poprawnie \\
    Wyniki obliczeń strat ciśnienia przedstawia tabela~3.4.  & poprawnie \\
    Wyniki obliczeń strat ciśnienia zestawiłem w~tabeli~3.4. & niepoprawnie \\
    W~naszej pracy policzyliśmy straty ciśnienia.            & niepoprawnie \\
    \bottomrule
  \end{tabularx}
\end{table}

Kilka zasad typograficznych, o~które warto zadbać od początku:
\begin{itemize}
  \item jednoliterowe spójniki i~przyimki (a, i, o, u, w, z) nie mogą
        zostawać na końcu wiersza -- po nich wpisuje się tyldę, np.\
        \verb|w~budynku|,
  \item liczbę od jednostki oddziela się spacją niełamliwą; najprościej
        użyć polecenia \verb|\qty|, np.\ \verb|\qty{2.5}{kW}| daje
        \qty{2.5}{kW} (z~przecinkiem dziesiętnym),
  \item zakresy liczbowe zapisuje się półpauzą: \verb|10--15| daje 10--15,
        a~w~połączeniu z~jednostką: \qtyrange{10}{15}{\degreeCelsius},
  \item cudzysłów polski ma postać „tekst” (\verb|,,tekst''|),
  \item skróty wyjaśnia się przy pierwszym użyciu, np.\ odnawialne źródła
        energii (OZE), i~zbiera w~wykazie oznaczeń.
\end{itemize}

\section{Wzory}
\label{sec:wzory}

Wzory wyrównuje się do środka, a~ich numery -- w~nawiasach okrągłych -- do
prawego marginesu. Numeracja jest prowadzona w~obrębie rozdziału. Wzór
stanowi część zdania, dlatego po nim stawia się przecinek lub kropkę.
Wielkości występujące we wzorze objaśnia się bezpośrednio pod nim
(środowisko \texttt{gdzie}), podając jednostki w~nawiasach kwadratowych.

\subsection{Wzór pojedynczy}

Współczynnik przenikania ciepła przegrody złożonej z~warstw jednorodnych
oblicza się zgodnie z~normą PN-EN ISO 6946 \cite{pn6946} ze wzoru
\begin{equation}
  U = \frac{1}{R_\mathrm{si} + \sum\limits_{i=1}^{n}\dfrac{d_i}{\lambda_i}
      + R_\mathrm{se}},
  \label{eq:U}
\end{equation}
\begin{gdzie}
  \zm{R_\mathrm{si},\ R_\mathrm{se}}{opór przejmowania ciepła na
      powierzchni wewnętrznej i~zewnętrznej}{\unit{m^2.K/W}}
  \zm{d_i}{grubość $i$-tej warstwy}{\unit{m}}
  \zm{\lambda_i}{obliczeniowy współczynnik przewodzenia ciepła materiału
      $i$-tej warstwy}{\unit{W/(m.K)}}
  \zm{n}{liczba warstw przegrody}{}
\end{gdzie}

Do wzoru odwołuje się poleceniem \verb|\eqref|, np.: \emph{wartość $U$
obliczona ze wzoru~\eqref{eq:U} dla ściany z~listingu~\ref{lst:U} wynosi
ok.\ \qty{0.23}{W/(m^2.K)}}. Zmienne zapisuje się kursywą, a~indeksy
opisowe (np.\ \emph{si} -- \emph{surface internal}) -- pismem prostym:
\verb|R_\mathrm{si}|. Dzięki pakietowi \texttt{icomma} przecinek
dziesiętny we wzorach nie tworzy zbędnego odstępu: $\theta_\mathrm{e} =
-20$~°C, $\eta = 0,92$.

\subsection{Układ kilku wzorów}

Kilka powiązanych zależności wyrównuje się względem znaku równości
(środowisko \texttt{align}). Spadek ciśnienia w~przewodzie opisuje wzór
Darcy'ego--Weisbacha~\eqref{eq:dw}, a~współczynnik oporów liniowych dla
przepływu burzliwego -- równanie Colebrooka--White'a~\eqref{eq:cw}:
\begin{align}
  \Delta p &= \lambda_\mathrm{f}\,\frac{L}{D}\,\frac{\rho\,w^2}{2},
  \label{eq:dw}\\
  \frac{1}{\sqrt{\lambda_\mathrm{f}}} &= -2\log_{10}\left(
    \frac{2,51}{\mathit{Re}\,\sqrt{\lambda_\mathrm{f}}}
    + \frac{k}{3,71\,D}\right),
  \label{eq:cw}
\end{align}
\begin{gdzie}
  \zm{\Delta p}{strata ciśnienia na długości przewodu}{\unit{Pa}}
  \zm{\lambda_\mathrm{f}}{współczynnik oporów liniowych}{}
  \zm{L,\ D}{długość i~średnica wewnętrzna przewodu}{\unit{m}}
  \zm{\rho}{gęstość wody}{\unit{kg/m^3}}
  \zm{w}{średnia prędkość przepływu}{\unit{m/s}}
  \zm{\mathit{Re}}{liczba Reynoldsa}{}
  \zm{k}{chropowatość bezwzględna ścianki}{\unit{m}}
\end{gdzie}

\section{Rysunki}
\label{sec:rysunki}

Rysunkiem jest każda ilustracja: wykres, schemat, mapa, zdjęcie czy zrzut
ekranu -- w~podpisie wszystkie oznacza się skrótem „Rys.”. Rysunek
wyśrodkowuje się, a~podpis (kursywą, bez kropki na końcu) umieszcza pod
nim. \textbf{Każdy rysunek musi zostać omówiony w~tekście}, a~odwołanie
powinno poprzedzać sam rysunek. Ilustracje rastrowe powinny mieć
rozdzielczość co najmniej \qty{300}{dpi}; wykresy najlepiej wstawiać jako
grafikę wektorową (PDF) albo rysować bezpośrednio w~\LaTeX-u. Jeżeli
rysunek pochodzi z~literatury lub powstał na podstawie cudzych danych,
w~podpisie podaje się źródło.

\subsection{Wykres rysowany w~\LaTeX-u}

Przykładową krzywą grzewczą, czyli zależność temperatury wody
zasilającej i~powrotnej od temperatury powietrza zewnętrznego, pokazano na
rys.~\ref{fig:krzywa}. Wykres powstał w~pakiecie \texttt{pgfplots}, więc
czcionki i~przecinki dziesiętne są spójne z~tekstem pracy.

\begin{figure}[htbp]
  \centering
  \begin{tikzpicture}
    \begin{axis}[
      width=0.82\linewidth, height=6.2cm,
      xmin=-20, xmax=20, ymin=15, ymax=75,
      x dir=reverse,
      xlabel={Temperatura zewnętrzna $\theta_\mathrm{e}$ [°C]},
      ylabel={Temperatura wody [°C]},
      grid=major, grid style={black!15},
      legend pos=north west, legend cell align=left,
      every axis plot/.append style={thick, samples=60, domain=-20:19},
    ]
      \addplot[black]        {20 + 50*((20 - x)/40)^0.77};
      \addplot[black, dashed]{20 + 35*((20 - x)/40)^0.77};
      \legend{zasilanie, powrót}
    \end{axis}
  \end{tikzpicture}
  \caption{Przykładowa krzywa grzewcza dla parametrów obliczeniowych
    70/55\,°C (dane umowne)}
  \label{fig:krzywa}
\end{figure}

\subsection{Schemat blokowy}

Proste schematy można rysować w~pakiecie TikZ. Rysunek~\ref{fig:schemat}
przedstawia uproszczony schemat systemu ciepłowniczego.

\begin{figure}[htbp]
  \centering
  \begin{tikzpicture}[
      node distance=8mm and 9mm,
      blok/.style={draw, thick, rounded corners=2pt, align=center,
                   minimum width=2.5cm, minimum height=1.2cm, font=\small},
      zas/.style={-{Latex[length=2.2mm]}, thick},
      pow/.style={-{Latex[length=2.2mm]}, thick, dashed}]
    \node[blok] (zrodlo) {Źródło\\ciepła};
    \node[blok, right=of zrodlo] (siec) {Sieć\\ciepłownicza};
    \node[blok, right=of siec]   (wezel) {Węzeł\\cieplny};
    \node[blok, right=of wezel]  (odb) {Instalacje\\c.o. i~c.w.u.};
    \draw[zas] ([yshift=3mm]zrodlo.east) -- ([yshift=3mm]siec.west);
    \draw[zas] ([yshift=3mm]siec.east)   -- ([yshift=3mm]wezel.west);
    \draw[zas] ([yshift=3mm]wezel.east)  -- ([yshift=3mm]odb.west);
    \draw[pow] ([yshift=-3mm]odb.west)   -- ([yshift=-3mm]wezel.east);
    \draw[pow] ([yshift=-3mm]wezel.west) -- ([yshift=-3mm]siec.east);
    \draw[pow] ([yshift=-3mm]siec.west)  -- ([yshift=-3mm]zrodlo.east);
  \end{tikzpicture}
  \caption{Schemat systemu ciepłowniczego (linia ciągła -- zasilanie,
    przerywana -- powrót)}
  \label{fig:schemat}
\end{figure}

\subsection{Rysunek z~pliku}

Do wstawiania plików graficznych służy polecenie \verb|\includegraphics|
(rys.~\ref{fig:plik}), a~same pliki umieszcza się w~katalogu
\texttt{rysunki/}. Rozszerzenie pliku można pominąć. Użyte tu obrazy zastępcze (pakiet \texttt{mwe}) są
dostępne w~każdej dystrybucji \LaTeX-a -- należy je zastąpić własnymi
ilustracjami.

\begin{figure}[htbp]
  \centering
  \includegraphics[width=0.55\linewidth]{example-image-a}
  \caption{Rysunek wstawiony z~pliku graficznego (źródło: opracowanie
    własne na podstawie \cite{pvgis})}
  \label{fig:plik}
\end{figure}

\subsection{Rysunek złożony}

Kilka powiązanych ilustracji można zestawić obok siebie jako jeden rysunek
z~podpisami cząstkowymi (rys.~\ref{fig:zlozony}). Do części rysunku
odwołuje się tak samo: rys.~\ref{fig:zlozony-a}.

\begin{figure}[htbp]
  \centering
  \begin{subfigure}[b]{0.45\linewidth}
    \centering
    \includegraphics[width=\linewidth]{example-image-b}
    \caption{stan przed modernizacją}
    \label{fig:zlozony-a}
  \end{subfigure}\hfill
  \begin{subfigure}[b]{0.45\linewidth}
    \centering
    \includegraphics[width=\linewidth]{example-image-c}
    \caption{stan po modernizacji}
    \label{fig:zlozony-b}
  \end{subfigure}
  \caption{Przykład rysunku złożonego z~dwóch części}
  \label{fig:zlozony}
\end{figure}

\section{Tabele}
\label{sec:tabele}

Tabelę poprzedza podpis umieszczony \emph{nad} nią. Tak jak rysunek, musi
ona zostać przywołana i~skomentowana w~tekście. Jednostki podaje się
w~nagłówkach kolumn, a~liczby w~kolumnie wyrównuje względem przecinka
dziesiętnego (kolumna typu \texttt{S} z~pakietu \texttt{siunitx}).

\subsection{Tabela wpisana ręcznie}

Orientacyjne wartości współczynnika przewodzenia ciepła wybranych
materiałów budowlanych zestawiono w~tabeli~\ref{tab:lambda}.

\begin{table}[htbp]
  \centering
  \caption{Współczynnik przewodzenia ciepła wybranych materiałów (wartości
    orientacyjne)}
  \label{tab:lambda}
  \begin{tabular}{@{}l S[table-format=4.0] S[table-format=1.3]@{}}
    \toprule
    \textbf{Materiał} & {\textbf{Gęstość} $\rho$ [\unit{kg/m^3}]}
                      & {$\lambda$ [\unit{W/(m.K)}]} \\
    \midrule
    beton zwykły                          & 2400 & 1.70  \\
    mur z~cegły ceramicznej pełnej        & 1800 & 0.77  \\
    tynk cementowo-wapienny               & 1850 & 0.82  \\
    drewno sosnowe (w~poprzek włókien)    & 550  & 0.16  \\
    wełna mineralna                       & 40   & 0.035 \\
    styropian EPS                         & 15   & 0.040 \\
    \bottomrule
  \end{tabular}
\end{table}

\subsection{Tabela wczytana z~pliku CSV}

Dane z~arkusza kalkulacyjnego nie muszą być przepisywane ręcznie. Wystarczy
zapisać je jako plik CSV (w~polskiej wersji Excela separatorem jest średnik,
a~znakiem dziesiętnym przecinek) i~wczytać pakietem \texttt{pgfplotstable}.
Tabela~\ref{tab:csv} powstała w~ten sposób z~pliku
\texttt{dane/zuzycie\_ciepla.csv}.

\begin{table}[htbp]
  \centering
  \caption{Miesięczne zużycie ciepła w~budynku przykładowym (dane umowne)}
  \label{tab:csv}
  \pgfplotstabletypeset[
    col sep=semicolon,
    /pgf/number format/read comma as period,
    columns/miesiac/.style={string type, column name={Miesiąc},
                            column type={c}},
    columns/t_zewn/.style={column name={$\theta_\mathrm{e}$ [°C]},
                           fixed, fixed zerofill, precision=1,
                           column type={r}},
    columns/Q_co/.style={column name={$Q_\text{c.o.}$ [GJ]},
                         fixed, fixed zerofill, precision=1,
                         column type={r}},
    columns/Q_cwu/.style={column name={$Q_\text{c.w.u.}$ [GJ]},
                          fixed, fixed zerofill, precision=1,
                          column type={r}},
    every head row/.style={before row=\toprule, after row=\midrule},
    every last row/.style={after row=\bottomrule},
  ]{dane/zuzycie_ciepla.csv}
\end{table}

Z~tego samego pliku można od razu sporządzić wykres
(rys.~\ref{fig:csv}) -- po zmianie danych w~pliku CSV zarówno tabela, jak
i~wykres zaktualizują się przy kolejnej kompilacji.

\pgfplotstableread[col sep=semicolon]{dane/zuzycie_ciepla.csv}\tabCieplo
\begin{figure}[htbp]
  \centering
  \pgfplotsset{/pgf/number format/read comma as period}
  \begin{tikzpicture}
    \begin{axis}[
      width=0.9\linewidth, height=6cm,
      ybar stacked, bar width=11pt, ymin=0,
      xtick=data, xticklabels from table={\tabCieplo}{miesiac},
      xlabel={Miesiąc}, ylabel={Zużycie ciepła [GJ]},
      ymajorgrids, grid style={black!15},
      legend style={at={(0.5,0.97)}, anchor=north, legend columns=2,
                    draw=none, font=\small},
      enlarge x limits=0.05,
    ]
      \addplot[fill=black!55, draw=black]
        table[x expr=\coordindex, y=Q_co]{\tabCieplo};
      \addplot[fill=black!15, draw=black]
        table[x expr=\coordindex, y=Q_cwu]{\tabCieplo};
      \legend{c.o., c.w.u.}
    \end{axis}
  \end{tikzpicture}
  \caption{Wykres sporządzony na podstawie danych z~tabeli~\ref{tab:csv}}
  \label{fig:csv}
\end{figure}

\section{Fragmenty kodu}
\label{sec:kod}

Krótkie fragmenty programów wstawia się w~środowisku \texttt{lstlisting}
(listing~\ref{lst:U}), dłuższe -- z~pliku poleceniem
\verb|\lstinputlisting| i~najlepiej w~załączniku (listing~\ref{lst:dp}
w~załączniku~\ref{zal:kod}). Polskie znaki w~komentarzach są obsługiwane.

\begin{lstlisting}[language=Python, caption={Obliczanie współczynnika
  przenikania ciepła przegrody}, label={lst:U}]
def wspolczynnik_U(warstwy, R_si=0.13, R_se=0.04):
    """Zwraca współczynnik przenikania ciepła U [W/(m2K)].

    warstwy - lista krotek (grubość d [m], lambda [W/(mK)])
    """
    R_t = R_si + R_se + sum(d / lam for d, lam in warstwy)
    return 1.0 / R_t

# Ściana: tynk 1,5 cm, cegła pełna 25 cm, EPS 15 cm, tynk 0,5 cm
sciana = [(0.015, 0.82), (0.25, 0.77), (0.15, 0.040), (0.005, 0.82)]
print(f"U = {wspolczynnik_U(sciana):.3f} W/(m2K)")
\end{lstlisting}

\section{Wyliczenia i~przypisy}

Wyliczenia punktowane stosuje się dla elementów równorzędnych, numerowane
-- gdy liczy się kolejność, np.\ dla etapów obliczeń:
\begin{enumerate}
  \item zebranie danych o~budynku i~lokalizacji,
  \item obliczenie strat ciepła przez przenikanie i~wentylację:
    \begin{enumerate}
      \item dla poszczególnych pomieszczeń,
      \item dla całego budynku,
    \end{enumerate}
  \item dobór źródła ciepła.
\end{enumerate}
Przypisy dolne służą do krótkich uzupełnień, które przerywałyby tok
wywodu\footnote{Przypisów nie należy używać zamiast wykazu literatury --
źródła podaje się w~nawiasach kwadratowych w~tekście.}.
